\section[Implementation in gretl]{Implementation in \pkg{gretl}}\label{BmaGretlSec3}
In this section, we describe the code as well as the use of the \pkg{gretl}
package for Bayesian model averaging, together with the accompanying
jointness measures. First, we will characterize our code and the use
of the GUI, then we will present how to use \pkg{BMA} via a
script. Finally, we will show the outputs that are returned.


\subsection[The core of the BMA code]{The core of the \pkg{BMA} code}
The \pkg{BMA} package consists of 3 public and 22 private functions, but
only 16 of them are used regularly. The name of each function starts
with the prefix \code{BMA_}.

\subsubsection{Public functions}
The public functions are:
\begin{enumerate}
	\item \code{BMA_GUI}
	\item \code{BMA_Print} 
	\item \code{BMA_OldStyle}
\end{enumerate}

The core package function which runs and controls the main loop is:
\begin{CodeChunk}
	\begin{CodeInput}
		function bundle BMA_GUI (series Y "Dependent variable",
		list X_list "List of indenpendent variables (X)",
		int acc_type[1:2:1] "Model prior" {"Binomial", "Binomial-Beta"},
		scalar av_model_size[0::1] "Prior average model size",
		scalar alpha[0:1:0.6] "Significance level for the initial model",
		int l_rank[2::4] "Number of the top ranked models",
		int g_type[1:5:1] "g-prior type" {"Benchmark prior",
			"Unit Information Prior (g-UIP)", "Risk Inflation Criterion (g-RIC)",
			"Hannan and Quinn HQC", "Root of g-UIP"},
		int do_joint[0:2:0] "Jointness analysis" {"None", "Ley-Steel Measure",
			"Doppelhofer-Weeks Measure"},
		int h_predict[0::0] "Number of out-of-sample forecasts",
		int Nrep[1e+06] "Total number of replications",
		int burn[0:99:10] "Percentage of burn-in draws",
		int verbosity[1:2:1] "Verbosity",
		bool show_progress[1] "Show overall progress",
		bool force_dof_lack[0] "Force more variables than observations")
	\end{CodeInput}
\end{CodeChunk}
%
The words in quotation marks are the labels for the GUI shown in
Figure~\ref{fig:Main window for BMA}.
\begin{itemize}
	\item The \code{series Y} is just a series of dependent variable.
	\item The \code{X_list} is a \pkg{gretl}'s object ``named list'' which
	is just a set of \textit{K} indenpendent variables (defined by names or dataset
	IDs). What is very important the \code{X_list} cannot contain a
	\code{const} (\pkg{gretl}'s internal and automatically generated
	constant term).
	\item The \code{acc_type[1:2:1]} is an integer indicating model prior
	type (the default is binomial, see
	Equations~\ref{eq:model_prior_binomial} and
	\ref{eq:model_prior_binomial-beta}).
	\item The \code{scalar av_model_size[0::1]} is the scalar of the prior
	average model size. Note, if \code{av_model_size} $ = \frac{K}{2}$
	and the model prior is set to binomial, we get the uniform prior on
	the model space. The default value is 1.
	\item The \code{scalar alpha[0:1:0.6]} is the significance level in
	OLS estimation. An independent variable enters the initial model if
	its $p$~value is less than the significance level (see
	\citealt{Koop2003}). The default value is $\alpha=0.6$, but setting
	$\alpha=1$ results in a model consisting of 0 to \textit{K} randomly
	chosen explanatory variables.
	\item The \code{int l_rank[2::10]} is the number of the top ranked
	models.  The default value is 10.
	\item The \code{int g_type[1:5:1]} indicates the type of g-prior to be
	used (the default value is 1: ``Benchmark prior''), see
	Equations~\ref{eq:g-UIP}, \ref{eq:g-RIC}, \ref{eq:Benchmark_Prior},
	\ref{eq:g-HQ} and~\ref{eq:root-g-UIP}.
	\item The \code{int do_joint[0:2:0]} indicates whether we do jointness
	analysis and if so, which measure to use. The default value is 0:
	``None'', see Section~\ref{sec:jointness_measures}.
	\item The \code{h_predict[0::]} is the number of out-of-sample
	forecasts.
	\item The \code{int Nrep[1000000]} is the total number of replications
	in the Monte Carlo simulation.  The default value is 1000000.
	\item The \code{int burn[0:99:10]} is the percentage of burn-in draws
	ranging from 0\% to 99\%. The default value is 10\%.
	\item The \code{int verbosity[1:2:1]} indicates the level of verbosity
	of the \pkg{BMA} package for printing results. The default value is 1,
	i.e., silent mode.
	\item The \code{bool show_progress[1]} switches progress printing during
	the simulation. Note that it consumes a lot of CPU time and significantly
	slows the package.
	\item The \code{bool force_dof_lack[0]} makes the BMA analysis to run even
	if there is not sufficient degrees of freedom.
\end{itemize}

Function for results printing is: 
\begin{CodeChunk}
	\begin{CodeInput}
		function void BMA_Print (bundle *MODEL)
	\end{CodeInput}
\end{CodeChunk}
This function takes as an argument reference to the bundle MODEL
with BMA results and prints results according to \code{verbosity} settings.

The last public function is \code{BMA_OldStyle} which is written
for backwards compatibility and can be used with scripts prepared
for old versions of BMA package (prior to 2.0).


\subsubsection{The main loop}
The main loop of the \pkg{BMA} package is split into five parts
controlled by following functions:
\begin{enumerate}
	\item \code{BMA_parse} -- basic parsing and dataset checks,
	\item \code{BMA_setup} -- bundle setup,
	\item \code{BMA_main} -- MCMC,
	\item \code{BMA_summary} -- final computations and bundle clean up,
	\item \code{BMA_Print} -- results printing.
\end{enumerate}

Essentials of BMA analysis is the MCMC stuff driven by \code{BMA_main}.
This part is split into thee parts:
\begin{enumerate}
	\item Setting up the $\mbox{MC}^3$ sampling algorithm.
	\item Constructing the initial model.
	\item $\mbox{MC}^3$ simulation.
\end{enumerate}


\subsubsection[The bundle returned by the BMA package]{The bundle returned by the \pkg{BMA} package}
The \pkg{BMA} package can optionally return bundle containing
results obtained in the analysis. The most interesting are:
\begin{itemize}
	\item \code{Post_Results} containing posterior moments
	(unconditional and conditional on inclusion),
	\item \code{jointness_m} containing posterior joint probability of variables,
	\item \code{prediction} containing predictive results.
\end{itemize}


The structure of the \code{Post_Results} matrix is
shown in Figure~\ref{fig:BMA_Post_Results_matrix_structure}.
It has \textit{K} rows (one for each explanatory variable)
and five columns are: the posterior inclusion probabilities, the
posterior mean and standard deviation of each coefficient (Mean and
Std.Dev.) and the posterior mean and standard deviation of each
coefficient conditional on the variable being included in the model
(Cond.Mean and Cond.Std.Dev.), see page~\pageref{sec:outputs} for
details.

\begin{figure}[ht!]
	\centering
	\begin{tabular}{cccccc}\hline
		& PIP & Mean & Std.Dev. & Cond.Mean & Cond.Std.Dev.\\
		$x_1$ & & & & &\\
		$x_2$ & & & & &\\
		$x_3$ & & & & &\\
		\vdots & & & & &\\
		$x_K$ & & & & &\\\hline
	\end{tabular}
	\caption{Structure of the \code{Post\_Results} matrix returned by the \pkg{BMA} package.}
	\label{fig:BMA_Post_Results_matrix_structure}
\end{figure}

The structure of the \code{jointness_m} matrix is
shown in Figure~\ref{fig:BMA_jointness_m_matrix_structure}.

\begin{figure}[ht!]
	\centering
	\begin{tabular}{cccccc}\hline
		& $x_1$ &$x_2$ &$x_3$ &\ldots&$x_K$\\
		$x_1$  &&&&&\\
		$x_2$  &&&&&\\
		$x_3$  & \multicolumn{5}{c}{joint probability of variables}\\
		\vdots &&&&&\\
		$x_K$  &&&&&\\\hline
	\end{tabular}
	\caption{Structure of the \code{jointness\_m} matrix returned by the \pkg{BMA} package.}
	\label{fig:BMA_jointness_m_matrix_structure}
\end{figure}
 
